\chapter{TSH (Thyroid-Stimulating Hormone)}

\section*{What Is This Test?}
Thyroid-stimulating hormone (TSH, also called thyrotropin) is a glycoprotein hormone produced by thyrotrope cells in the anterior pituitary gland. It stimulates the thyroid gland to produce and release thyroxine (T4) and triiodothyronine (T3). TSH is regulated by thyrotropin-releasing hormone (TRH) from the hypothalamus and negative feedback from circulating T3 and T4. The TSH test measures the serum concentration of thyrotropin using high-sensitivity immunometric assays (third-generation assays) capable of detecting TSH at concentrations below 0.01 mIU/L. TSH is the single best screening test for thyroid dysfunction because of the logarithmic relationship between TSH and free T4 --- small changes in free T4 produce large changes in TSH. Current third-generation assays have a functional sensitivity of 0.01--0.02 mIU/L, enabling reliable discrimination between normal and suppressed TSH levels.

\section*{Alternative Names}
Thyrotropin; Thyroid-Stimulating Hormone; Thyrotropic Hormone; TSH; TSH Third Generation; High-Sensitivity TSH; hTSH

\section*{Principle}
TSH is measured using a two-site sandwich immunometric assay (noncompetitive immunoassay) on automated analyzers. One antibody (capture antibody) is bound to a solid phase (magnetic particles or microtiter wells) and recognizes one epitope on the TSH molecule. A second antibody (detection antibody) is labeled with a chemiluminescent tag (e.g., ruthenium complex in Roche ECLIA, acridinium ester in Siemens CLIA, alkaline phosphatase in Beckman Coulter) and binds a different epitope. After incubation, unbound material is washed away and the chemiluminescent signal is proportional to TSH concentration. Most platforms use electrochemiluminescence immunoassay (ECLIA) on Roche Cobas analyzers, chemiluminescent microparticle immunoassay (CMIA) on Abbott Architect, or chemiluminescent immunoassay (CLIA) on Siemens Atellica/ADVIA Centaur. Third-generation assays use chemiluminescent detection with functional sensitivity (lowest concentration measurable with a coefficient of variation <20\%) at or below 0.01--0.02 mIU/L. Calibration is standardized against the WHO 2nd International Reference Preparation (IRP 80/558).

\section*{History}
TSH was first isolated from bovine pituitary glands in the 1930s. The first radioimmunoassay (RIA) for TSH was developed by Utiger and Odell in 1965, with a detection limit of approximately 1.0 mIU/L, which could only detect elevated TSH levels. The 1980s saw development of first-generation immunometric assays (IRMA) with sensitivity around 0.1 mIU/L. Second-generation assays (1990s) used chemiluminescent labels with sensitivity of 0.01--0.02 mIU/L. Third-generation assays (2000s to present) achieved functional sensitivity <0.01 mIU/L, enabling precise classification of subclinical hyperthyroidism. The current trend is toward harmonization of TSH assays across manufacturers using shared reference standards, though significant inter-assay variability persists (approximately 10--15\% CV between platforms).

\section*{Why This Test Is Ordered}
\begin{itemize}
  \item Primary screening test for thyroid dysfunction in symptomatic patients (fatigue, weight changes, temperature intolerance, palpitations, neck swelling) or during routine health examinations
  \item Diagnosis of primary hypothyroidism (elevated TSH >4.5 mIU/L) or primary hyperthyroidism (suppressed TSH <0.3--0.5 mIU/L depending on assay)
  \item Monitoring levothyroxine replacement therapy dose adequacy in hypothyroid patients, with target TSH typically 0.5--2.5 mIU/L
  \item Newborn screening for congenital hypothyroidism (elevated TSH >20 mIU/L in heel-stick blood spot) to enable early treatment and prevent intellectual disability
  \item Evaluation of subclinical thyroid dysfunction (TSH 4.5--10 mIU/L with normal free T4) in patients with cardiovascular risk factors including atrial fibrillation, osteoporosis, or hyperlipidemia
  \item Monitoring thyroid hormone suppression therapy in differentiated thyroid cancer patients (target TSH <0.1--0.5 mIU/L)
  \item Assessment of pituitary function in suspected hypopituitarism or postpartum Sheehan syndrome
  \item Serial monitoring of TSH in amiodarone-treated patients (10--15\% develop thyroid dysfunction) at baseline and every 6 months
  \item Screening for thyroid dysfunction during pregnancy (first trimester TSH upper limit 2.5 mIU/L)
  \item Investigation of goiter or thyroid nodules
\end{itemize}

\section*{Primary vs Secondary Test}
This is a primary screening test for thyroid dysfunction. TSH is the first-line and most sensitive test for evaluating thyroid status, recommended by the American Thyroid Association, Endocrine Society, and American Association of Clinical Endocrinologists as the initial test for suspected thyroid disease in outpatients. Free T4 and/or free T3 are reflexively added when TSH is abnormal.

\section*{How This Test Is Performed}
A healthcare professional draws a blood sample from a vein in your arm using a small needle into a serum separator tube (SST, gold-top or tiger-top tube) or plain red-top tube. After the needle is inserted, approximately 3--5 mL of blood is collected. You may feel a slight sting when the needle is inserted or removed. The procedure typically takes less than 5 minutes. The specimen is centrifuged to separate serum, which is stable at 2--8 degrees C for 7 days or at -20 degrees C for 6 months. TSH exhibits diurnal variation with a nadir in late afternoon and peak around midnight, so samples should ideally be drawn at the same time of day for serial monitoring.

\section*{What Preparation Is Needed}
No special preparation is required. TSH testing does not require fasting, though morning collection (before 10 AM) is preferred due to circadian variation. Patients taking biotin (vitamin B7) supplements at high doses (>5 mg/day) should discontinue biotin for at least 24--48 hours before testing because biotin interferes with streptavidin-biotin immunoassay platforms (Roche, Siemens, Beckman Coulter), potentially causing falsely low TSH results. Patients on levothyroxine should take their blood sample before their daily dose (trough level) for consistent monitoring. For pregnancy screening, testing should occur at the first prenatal visit.

\section*{Contraindications and Precautions}
No absolute contraindications. Factors affecting TSH results include: acute non-thyroidal illness (euthyroid sick syndrome can transiently suppress TSH); dopamine, dobutamine, and high-dose glucocorticoids (suppress TSH); amiodarone (elevates or suppresses TSH depending on phase); lithium (elevates TSH in 20--30\% of users); phenytoin, carbamazepine, and rifampin (lower free T4 and may elevate TSH); recent radioiodine contrast agents (iodine load affects thyroid function for 2--4 weeks); pregnancy (hCG cross-reacts with TSH receptor, physiologically suppressing TSH in first trimester); heterophile antibodies (including human anti-mouse antibodies) can interfere with sandwich immunoassays causing falsely elevated TSH. Specimen hemolysis or lipemia can interfere with optical detection methods. TSH is more stable than free T4/T3 --- samples can be stored refrigerated for 7 days without significant degradation.

\section*{Risks}
Minimal risk. Slight pain, bruising, or hematoma at venipuncture site. Rarely: vasovagal syncope, infection at puncture site, or phlebitis. Excessive bleeding risk is minimal in patients with normal coagulation.

\section*{What Happens After the Test}
Results are typically available within 1--4 hours on modern automated platforms. The healthcare provider interprets TSH in conjunction with clinical symptoms, physical examination findings, and often reflex free T4/T3 testing. If TSH is abnormal and untreated, follow-up typically occurs within 1--4 weeks to initiate or adjust therapy. Stable treated patients are monitored every 6--12 months. Abnormal results may prompt: free T4 and free T3 measurement, thyroid antibodies (anti-TPO, anti-thyroglobulin, TSH-receptor antibodies), thyroid ultrasound, radioactive iodine uptake scan, or endocrinology referral.

\section*{Understanding Results}

\subsection*{Below Normal}
\begin{itemize}
  \item \textbf{Suppressed TSH (<0.3 mIU/L):} Primary hyperthyroidism --- most commonly Graves disease (autoimmune, thyroid-stimulating immunoglobulin activation of TSH receptor) producing diffuse toxic goiter, or toxic multinodular goiter with autonomous hyperfunctioning nodules.
  \item \textbf{Subclinical hyperthyroidism (TSH 0.1--0.4 mIU/L, normal free T4):} Associated with atrial fibrillation (3-fold increased risk), osteoporosis in postmenopausal women, and progression to overt hyperthyroidism at 1--5\% per year.
  \item \textbf{TSH-secreting pituitary adenoma (thyrotropinoma):} Rare cause (<1\% of pituitary adenomas) with inappropriately normal or elevated TSH despite elevated T4/T3 --- requires pituitary MRI and alpha-subunit measurement.
  \item \textbf{Exogenous thyrotoxicosis:} Surreptitious or iatrogenic levothyroxine over-replacement, most commonly from weight-loss preparations containing thyroid hormone.
  \item \textbf{First trimester of pregnancy:} Physiologic TSH suppression from hCG-mediated thyroid stimulation (resolves by second trimester).
  \item \textbf{Non-thyroidal illness (euthyroid sick syndrome):} Transient TSH suppression in critically ill patients receiving dopamine or high-dose glucocorticoids.
  \item \textbf{Subacute thyroiditis (de Quervain):} Initial thyrotoxic phase with suppressed TSH and elevated T4/T3 from follicular destruction, followed by hypothyroid phase.
\end{itemize}

\subsection*{Above Normal}
\begin{itemize}
  \item \textbf{Elevated TSH (>4.5 mIU/L) with low free T4:} Primary hypothyroidism --- most commonly Hashimoto thyroiditis (autoimmune lymphocytic infiltration causing progressive follicular destruction), detectable by elevated anti-TPO antibodies in >90\% of cases.
  \item \textbf{Subclinical hypothyroidism (TSH 4.5--10 mIU/L, normal free T4):} Affects 4--10\% of the general population, more common in women and with advancing age. Associated with increased cardiovascular risk, dyslipidemia, and a 2--5\% per year progression rate to overt hypothyroidism.
  \item \textbf{Iatrogenic hypothyroidism:} Post-thyroidectomy (total or subtotal), radioactive iodine (I-131) therapy for Graves disease or thyroid cancer, external beam radiation to neck.
  \item \textbf{Drug-induced:} Amiodarone (3--15\% develop hypothyroidism, more common in iodine-sufficient regions), lithium (inhibits thyroid hormone release, 20--30\% develop goiter or hypothyroidism), tyrosine kinase inhibitors (sunitinib, sorafenib).
  \item \textbf{TSH resistance due to inactivating TSH receptor mutations:} Rare congenital condition with elevated TSH, normal or low free T4, and hypoplastic thyroid gland --- diagnosed in infancy with normal newborn screening TSH.
  \item \textbf{Recovery phase of non-thyroidal illness:} Transient TSH elevation during recovery from severe illness.
  \item \textbf{Thyroid hormone resistance (THR-beta mutations):} Elevated TSH with elevated free T4/T3 --- autosomal dominant with variable phenotype from asymptomatic to goiter and ADHD. Prevalence approximately 1 in 40,000.
  \item \textbf{Pituitary hyperplasia from long-standing primary hypothyroidism:} Can appear as pituitary mass on MRI, resolves with levothyroxine treatment.
\end{itemize}

\section*{Normal Results}
\begin{itemize}
  \item \textbf{Adults (18--60 years):} 0.5--4.5 mIU/L (may vary by 0.3--5.0 mIU/L depending on laboratory reference)
  \item \textbf{Elderly (>60 years):} Upper limit may extend to 6.0--8.0 mIU/L (age-related increase in TSH set point)
  \item \textbf{First trimester pregnancy:} 0.1--2.5 mIU/L (lower limit due to hCG cross-reactivity)
  \item \textbf{Second trimester pregnancy:} 0.2--3.0 mIU/L
  \item \textbf{Third trimester pregnancy:} 0.3--3.0 mIU/L
  \item \textbf{Neonates (1--4 days):} 1.0--17.0 mIU/L (physiological TSH surge at birth, peaks at 30 minutes, declines over first week)
  \item \textbf{Children (1--18 years):} 0.5--5.0 mIU/L
  \item \textbf{Target range on levothyroxine therapy:} 0.5--2.5 mIU/L
  \item \textbf{Target range in thyroid cancer suppression:} <0.1 mIU/L (high-risk) to <0.5 mIU/L (low-risk)
\end{itemize}

\section*{Follow-up Tests}
\begin{itemize}
  \item \textbf{Free T4 (FT4):} Reflex testing in any patient with abnormal TSH. Assesses circulating unbound T4, which represents the biologically active hormone available to tissues. Used to determine severity of thyroid dysfunction and monitor treatment response.
  \item \textbf{Free T3 (FT3):} Added when TSH is suppressed with normal FT4 to evaluate for T3 thyrotoxicosis. Particularly useful in Graves disease and toxic nodular disease where T3 is disproportionately elevated.
  \item \textbf{Anti-thyroid peroxidase (anti-TPO) antibodies:} Confirm autoimmune etiology of hypothyroidism (Hashimoto thyroiditis). Elevated in >90\% of Hashimoto and 50--80\% of Graves disease.
  \item \textbf{Anti-thyroglobulin antibodies:} Present in 60--80\% of Hashimoto thyroiditis and 30--50\% of Graves disease. Also measured alongside thyroglobulin as a tumor marker in differentiated thyroid cancer follow-up.
  \item \textbf{TSH receptor antibodies (TRAb) / thyroid-stimulating immunoglobulins (TSI):} Specific for Graves disease. TRAb measured in pregnancy to predict risk of neonatal Graves disease (transplacental antibody transfer).
  \item \textbf{Thyroid ultrasound:} Imaging for thyroid nodules >1 cm, goiter, or suspicion of structural abnormalities. Doppler flow patterns differentiate Graves disease (hypervascular) from thyroiditis.
  \item \textbf{Radioactive iodine uptake (RAIU) and scan (I-123):} Distinguishes Graves disease (diffusely increased uptake, 35--80\% at 24 hours) from subacute thyroiditis (low/absent uptake, <5\%) and toxic nodular disease (focal uptake).
  \item \textbf{Pituitary MRI with contrast:} If TSH-secreting pituitary adenoma is suspected (inappropriately normal/elevated TSH with elevated free T4/T3).
  \item \textbf{Thyroglobulin:} Used as tumor marker in differentiated thyroid cancer follow-up after total thyroidectomy and ablation.
\end{itemize}

\section*{References}
\begin{enumerate}
  \item Garber JR, Cobin RH, Gharib H, et al. Clinical practice guidelines for hypothyroidism in adults: cosponsored by the American Association of Clinical Endocrinologists and the American Thyroid Association. Thyroid. 2012;22(12):1200-1235.
  \item Ross DS, Burch HB, Cooper DS, et al. 2016 American Thyroid Association guidelines for diagnosis and management of hyperthyroidism and other causes of thyrotoxicosis. Thyroid. 2016;26(10):1343-1421.
  \item Jonklaas J, Bianco AC, Bauer AJ, et al. Guidelines for the treatment of hypothyroidism: prepared by the American Thyroid Association Task Force on Thyroid Hormone Replacement. Thyroid. 2014;24(12):1670-1751.
  \item Alexander EK, Pearce EN, Brent GA, et al. 2017 Guidelines of the American Thyroid Association for the diagnosis and management of thyroid disease during pregnancy and the postpartum. Thyroid. 2017;27(3):315-389.
  \item Baloch Z, Carayon P, Conte-Devolx B, et al. Laboratory medicine practice guidelines: laboratory support for the diagnosis and monitoring of thyroid disease. Thyroid. 2003;13(1):3-126.
  \item Burtis CA, Ashwood ER, Bruns DE. Tietz Textbook of Clinical Chemistry and Molecular Diagnostics. 6th ed. Elsevier; 2023.
\end{enumerate}

\clearpage
\section*{Regulatory Status and Authorization Date}
Regulatory status applies to the specific assay, instrument, manufacturer, and intended use rather than to the test name alone. No single approval date can be assigned to this test category. For a commercial assay, verify the current product in the relevant regulator's database: the U.S. Food and Drug Administration (FDA) clearance, approval, or authorization; India's Central Drugs Standard Control Organisation (CDSCO) licence or permission; the European Union In Vitro Diagnostic Regulation (EU IVDR) CE certificate issued through the applicable conformity-assessment route; or the United Kingdom Medicines and Healthcare products Regulatory Agency (MHRA) registration and UKCA/CE status. Laboratory-developed tests may not have an individual FDA, CDSCO, EU IVDR, or MHRA approval date.
\clearpage
